% Product 트랙 Day 8 슬라이드 본문 — 손으로 쓴 정본(한 프레임 = 한 쪽). Day7_슬라이드_body.tex 와 같은 형식.
% 래퍼: Day8_슬라이드.tex(슬라이드판) / Day8_슬라이드_노트.tex(노트판, 우측에 \note).
% 화면에는 결론만, 설명·근거·예외는 각 프레임의 \note{}에. 그림은 ../그림/Product_Day8/*.pdf.
% 모든 숫자는 Day8 워크북 완성 예시본(정본)과 day8_enterprise.py 에서. 시점 1 D+800 (2028-05-10) / 시점 2 D+1000 (2028-11-26, 봉인).

\newcolumntype{L}[1]{>{\raggedright\arraybackslash}p{#1}}
\newcommand{\ra}{$\rightarrow$}
\newcommand{\til}{\textasciitilde}
\newcommand{\keymsg}[1]{\par\vfill{\color{mDarkTeal}\bfseries #1}\par}
\newcommand{\figcap}[1]{\par\smallskip{\footnotesize\color{black!60}#1}\par}
\newcommand{\fig}[2][0.66]{\centering\includegraphics[height=#1\textheight,width=\textwidth,keepaspectratio]{#2}\par}
\newcommand{\subhead}[1]{{\color{mDarkTeal}\bfseries #1}\par}
\newenvironment{tbl}[2][\small]{\begin{center}\usebeamercolor[fg]{normal text}\setlength{\tabcolsep}{4pt}\renewcommand{\arraystretch}{1.15}#1\begin{tabular}{#2}\toprule}{\bottomrule\end{tabular}\end{center}}

% =====================================================================
% 도입
% =====================================================================
\begin{frame}{Product 트랙 Day 8 — 시리즈A와 첫 엔터프라이즈: 제품 회사로 남을 수 있는가}
\begin{itemize}
\item 8일 과정의 마지막 날 · Product 트랙 4일차 · 시점이 \textbf{둘} — \textbf{2028-05-10 (D+800)}: 온담 요구 47건 확정·텀시트 2안 / \textbf{2028-11-26 (D+1000)}: 에필로그, 봉인
\item 회사 상태(시점 1): \textbf{27명}(개발 14) · 유료 \textbf{58개사} · ARR 11.9억(온담 제외) · 현금 7.52억 / 8.9개월 · 그리고 21억짜리 계약서와 요구 47건
\item 오늘의 문서 5종 ㊱\til{}㊵ — 딜 데스크 · 경계·SLA·배상 · IR·재무모델 · 텀시트 레드팀 · 편익 원장 공동 확정본(= PM ⑳)
\end{itemize}
\keymsg{오늘의 네 질문 — 같은 21억이 두 회사에게 어떻게 다른 무게인가 / 47건을 어떻게 판정하는가 / 같은 문서를 반대편에서 쓸 때 무엇이 달라지는가 / 표준 문구의 값은 얼마인가}
\note{표지. 어제 D+600에 쓴 다섯 문서(㉛\til{}㉟) 중 ㉝(표기 규칙)·㉞(항목 4·8)·㉟(CR-044)와 PM Day4 워크북 ⑯·⑰·⑳을 손 닿는 곳에. 사전 읽기 정본 §3.4·§4.10\til{}4.13·§5·§7.4·§7.5·§9(약 25분). 시점 2(D+1000)는 워크북 §7에 봉인 — 학생이 정본 §1.5를 읽었더라도 판정은 마지막 30분. 예상 소요 2분.}
\end{frame}

\begin{frame}[shrink=10]{오늘의 흐름}
\begin{columns}[T]
\begin{column}{0.46\textwidth}
\subhead{오전 — 딜의 구조·5차원·경계 (제1\til{}3장)}
\begin{itemize}
\item \textbf{1교시} 90분 · 두 시점 · 21억의 해부 · 네 비율 · BATNA · 5차원 표 · 3분류 · 용량 잠식 · 계약 문구 · 즉시 탈락 셋 — 제1·2장
\item \textbf{2교시} 90분 · \textbf{실습 ①-a} ㊱ 딜 데스크(45분) + 경계 3영역 · SLA 3층(공급자 쪽) · AI 판정 오류 — 3.1\til{}3.3절
\item \textbf{3교시} 90분 · 배상 4구간 · 218문항 · 이관 — 3.4\til{}3.6절 + \textbf{실습 ①-b} ㊲(35분)
\end{itemize}
\end{column}
\begin{column}{0.52\textwidth}
\subhead{오후 — 자본·편익 원장·8일의 끝}
\begin{itemize}
\item \textbf{4교시} 90분 · IR 두 값 · 3표기 · 3케이스 · IC 8문 · 캡테이블 · 옵션풀 1.29\%p · 레드팀·2×2 — 제4장 + \textbf{실습 ②} ㊳(25분) · ㊴(25분)
\item \textbf{5교시} 75분 · ⑳를 반대편에서 · 세 서명 · 인수 기준의 모집단 · 두 시계 · 사슬 — 제5장 + \textbf{실습 ③} ㊵(25분)
\item \textbf{6교시} 75분 · \textbf{실습 ④} 8일 사슬 역추적(조별) · 발표
\item \textbf{마무리} 30분 · 에필로그 봉인 해제 · 8일 최종 명제 · 학술 배경 · Product 트랙을 마치며
\end{itemize}
\end{column}
\end{columns}
\keymsg{같은 문서를 반대편에서 쓴다 — Day4에 보낸 SAC·요구한 배상·서명받은 편익 원장을 오늘 받고, 막고, 서명한다}
\note{6교시 틀(90·90·90·90·75·75 + 마무리 30 = 540). ★ 시간 — ㊱ 45 / ㊲ 35 / ㊳ 25 / ㊴ 25 / ㊵ 25. 두 시점 운영: 모든 카드는 시점 1 값만; ㊱ 항목 5 마지막 줄 "에필로그 예측"을 학생이 적은 것을 확인한 뒤 마무리에서 §7을 연다. 인쇄용 빈 템플릿(00\_Day8\_템플릿팩\_빈양식.pdf 23쪽) 배부. 예상 소요 3분.}
\end{frame}

\begin{frame}{학습 목표 — 오늘이 끝나면}
\begin{columns}[T]
\begin{column}{0.5\textwidth}
\subhead{오전}
\begin{enumerate}
\item 계약 총액을 세 줄(일시 / 반복 / 통과)로 나누고, ARR인 줄을 가려낸다 — 21억 중 3.78
\item 요구 목록을 항목이 아니라 \textbf{차원}으로 먼저 판정하고, 3분류와 용량 잠식(82.1\%)을 계약 조건(상한 6 FTE)으로 바꾼다
\item SLA 3층을 공급자 쪽에서 그리고, AI 판정 오류를 \textbf{모집단 한정}으로 SLA에 넣되 배상에서 뺀다
\end{enumerate}
\end{column}
\begin{column}{0.5\textwidth}
\subhead{오후}
\begin{enumerate}
\setcounter{enumi}{3}
\item 배상 상한의 물리적 하한을 검산값(6.12)으로 계산하고 4구간을 조건으로 가른다
\item IR 덱에 두 값을 병기하고, 옵션풀 "표준 문구"의 값(1.29\%p = 25.4억)을 주식수로 계산한다
\item ⑳의 표에 output 열을 채우고, 인수 기준의 \textbf{모집단}을 쓴 뒤에만 서명한다
\end{enumerate}
\end{column}
\end{columns}
\keymsg{8일의 마지막 문장 — FN 2.4\%는 자동차 3세트 60건의 값이다. 온담의 4번째 세트에서는 한 번도 측정되지 않았다}
\note{교재 머리말·각 장 핵심 정리. 여섯 목표가 ㊱\til{}㊵ 다섯 문서와 1:1은 아니다 — 1·2가 ㊱, 3·4가 ㊲, 5가 ㊳·㊴, 6이 ㊵. 예상 소요 2분.}
\end{frame}

\begin{frame}{Day7에서 가져오는 것 — 여섯 줄}
\begin{tbl}[\footnotesize]{L{40mm}L{28mm}L{62mm}}
Day7에서 & 오늘 어디에 & 무엇이 되는가 \\ \midrule
㉝ 표기 규칙(4.9 vs 3.8 · 61 vs 49.8) & ㊳ 항목 2 & ARR \textbf{18.9 vs 15.68} — 규칙이 덱에 절반만 적용됐다 \\
㉞ 항목 4 — 162만 선례·중간안 & ㊲ 항목 5 & 배상 상한 \textbf{10\% vs 100\%} — 4구간 \\
㉞ 항목 8 — 재발 방지 6건 증빙 & ㊲ 항목 6 & 218문항 중 \textbf{제218문항}의 답 \\
㉞ 항목 5 — "자동차 3세트 60건, 온담 미측정" & ㊵ 항목 4 & \textbf{인수 기준의 모집단} — 최종 명제 \\
㉟ 항목 4·5 — Won't·CR-044 보류 & ㊱ 항목 3 & C-31 온프렘 — Won't가 요구로 돌아온다 \\
㉝ 항목 5 — 런웨이 11, 브릿지 필요 & ㊴ 항목 2 & SAFE 8억 \ra{} K-03 \ra{} 텀시트 2안 \\
\end{tbl}
\keymsg{그리고 PM 트랙 ⑯·⑰·⑳ — 보내던·요구하던·서명받던 쪽에서, 받는·막는·서명하는 쪽으로}
\note{교재 서문 표. 세 가지 미리 말하기 — 두 시점을 섞지 않는다 / 반대편이다 / 분모의 마지막 판(ARR의 분자, FN의 모집단). 예상 소요 3분.}
\end{frame}

% =====================================================================
\section{1교시 (90분) — 엔터프라이즈 딜의 구조 · 코드체계 5차원 · 3분류 · 용량 잠식}
% =====================================================================
\begin{frame}{이 교시에서 배우는 것}
\begin{itemize}
\item \textbf{두 시점} — D+800에 쓰고 D+1000은 봉인. 섞으면 criteria drift(Day7 ㉞ 항목 6)
\item \textbf{21억의 해부} — 구축 3.00(실소요 6.12) + 라이선스 11.34(연 3.78 = ARR) + 운영 6.66(원가 통과). 라인 70 중 30은 정책 밖
\item \textbf{네 비율} — 4.906(온담) / 37.0(IR, 과대) / 24.1(검산) / 82.1(용량, 덱에 없다) · \textbf{BATNA} 대한기공 2.4 × 3
\item \textbf{5차원 표} — 47건을 판정하기 전에 · \textbf{3분류} 제품화 27 / 유상 59 / 거절 52 mm · \textbf{용량} 138 = 14명 9.9개월 · \textbf{즉시 탈락 셋}
\end{itemize}
\keymsg{집중 리스크를 ARR 비중으로만 말하는 IR 덱은 거짓말을 하고 있다 — 정본 §5.6}
\note{교재 제1·2장. 1교시는 이론만 90분 — 2교시 ㊱ 실습의 입력이 전부 여기서 나온다(5차원 표의 1차 판정 열은 학생이 쓴다 — 강사는 네 열만 보이고 다섯째 열을 비운 채 넘긴다). 예상 소요 2분.}
\end{frame}

\begin{frame}{1.1 D+600 \ra{} D+1000 — 두 시점}
\fig[0.52]{fig_1_1_timeline.pdf}
\figcap{그림 1-1. 시점 1(D+800)에 47건 확정·텀시트 2안·대한기공 제안이 모여 있다. 시점 2(D+1000)는 봉인. RFI는 정본 표기 D+640이나 날짜 2027-11-02 = D+610.}
\keymsg{결정은 요구가 확정되는 날 이전에 끝나야 한다 — 서명일에는 이미 47건이고 선행 조건이 붙어 있다}
\note{교재 1.1절. 사건 — D+610 RFI+218문항 · D+617 회신(131/62/25) · D+640 제안서(커버리지 95\% 가정) · D+660 SAFE · D+690 실사 61.4\% → 6.12 · D+718 1차 31건 · D+760 온담 G5(71\% 경고) · D+780 대한기공 · D+787 한강 · D+795 코너스톤 · D+800 47건 확정. 정본 §7.1의 'D+900' 열 값을 시점 1에 쓴다. 예상 소요 4분.}
\end{frame}

\begin{frame}{1.2 21억의 해부 — ARR인 줄은 하나뿐이다}
\fig[0.52]{fig_1_2_deal_structure.pdf}
\figcap{그림 1-2. 구축 3.00(실소요 6.12) · 라이선스 11.34(연 3.78 = 라인 70 × 45만 × 12) · 운영 6.66(월 1,850만). 61점 30라인은 "2개점 = 1라인"의 임시 합의.}
\keymsg{IR 덱 18.9 = 11.9 + 7.0(청구 전액) · 검산 15.68 = 11.9 + 3.78(라이선스만) — 3.00으로 6.12어치를 약속하면 1년차 $-$4.04가 $-$7.16이 된다}
\note{교재 1.2절 · 정본 §5.2·E-01\til{}05. 두 법인(홀딩스 54라인 / 푸드시스템 16). 실소요 분해(2.35 + 0.49 + 0.60 + 0.48 + 1.20 + 1.00)와 운영 분해(400 + 467 + 333 + 150 + 500)는 [추가 설정]. Day6 ㉘ 밸류메트릭이 점포에서 무너지는 물증 — 70 중 30의 산식이 없다. 예상 소요 5분.}
\end{frame}

\begin{frame}{1.3 같은 21억의 네 비율 — 4.9 / 37.0 / 24.1 / 82.1}
\fig[0.58]{fig_1_3_asymmetry.pdf}
\figcap{그림 1-3. 온담에게 21/428 = 4.906\%(원가 허용범위 10.08억의 두 배). 딥게이지 IR 7.0/18.9 = 37.0(과대) · 검산 3.78/15.68 = 24.1 · 개발 용량 138/168 = 82.1(덱에 없다).}
\keymsg{37은 과대, 24.1은 적정, 82.1이 실제 리스크 — 계약이 깨지는 리스크가 아니라 성사되는 리스크}
\note{교재 1.3절 "작은 숫자로 먼저". 순SaaS 분모면 26.8\%(심화). 온담 쪽 4.906\%를 알아야 CFO 3주·218문항이 이해된다. 예상 소요 5분.}
\end{frame}

\begin{frame}{1.4 BATNA — 대한기공, 그리고 "자동차를 벗어날 것인가"}
\fig[0.55]{fig_1_4_batna.pdf}
\figcap{그림 1-4. 연 7.0 vs 2.4 · ARR 3.78 vs 2.40 · 공수 138(86) vs 0 · 확장 SAM vs 초기 SAM · 온담 거절 = 코너스톤 소멸.}
\keymsg{BATNA는 서명돼야 카드다 — 병행 체결. 답: 벗어나지 않는다, 온담은 측정 한 건이다}
\note{교재 1.4절 · 정본 §5.8·§8.6·E-15. Fisher \& Ury "develop your BATNA". 표면의 질문(21억) vs 실제 질문(초기 SAM 2,150 → 확장 11,800). ㊱ 항목 8 조건부 Go 일곱 조건은 2교시 실습에서. 예상 소요 4분.}
\end{frame}

\begin{frame}{2.1 47건을 하나씩 판정하면 6시간 — 5차원 표 먼저}
\begin{itemize}
\item 요구 목록은 저자를 따라 읽힌다 — 스폰서의 19건은 중요해 보이고 영업본부장의 5건은 부차적으로 보인다
\item C-05(마스터 구축, 나영선)와 C-36(ERP 연동 API, 박세연)은 같은 차원(마스터 없음) — 하나씩 판정하면 하나는 제품화, 하나는 거절이 된다
\item 우리 제품과 이 고객의 코드체계가 \textbf{갈리는 차원}을 먼저 판정하면, 47건은 그 판정을 옮겨 적는 별지가 된다
\item 저자별 19 / 11 / 9 / 5 / 3건 · 62 / 31 / 28 / 12 / 5 mm = 47건 138 mm(E-10). 5차원 밖(인프라·보안·운영) 20건 43 mm도 표가 있어야 보인다
\end{itemize}
\keymsg{6시간 → 40분 — 판정이 저자를 따라가지 않게 하는 유일한 방법이 차원 표다}
\note{교재 2.1절 · 정본 §5.4. 정본 §8.2 — CPO 윤하경이 "요구사항 정리자"가 되는지 시험받는 지점. 태산정밀 62항목(Day6 B-20) 전례. 예상 소요 4분.}
\end{frame}

\begin{frame}{2.2 다섯 차원 읽기 — 그리고 1차 판정}
\fig[0.52]{fig_2_1_five_dims.pdf}
\figcap{그림 2-1. 정본 §5.4 네 열 + 1차 판정. ② 마스터 없음은 거절, ④ 법정 기준은 ㊲ 배상 분리와 연동 — 둘이 즉시 탈락 조건.}
\keymsg{①은 재설계, ②는 매핑 대상 부재, ③은 1차 키, ④는 규제 리스크, ⑤는 Flow 절반 — 성격이 다르면 판정도 다르다}
\note{교재 2.2절. 학생용 슬라이드에서는 다섯째 열을 가린 판을 먼저 보이고(강사 재량 — 손으로 가려도 된다), 2교시 실습 후 이 판을 보인다. ② 거절의 이유 둘 — 판정 근거의 소유 역전, 3년 유지보수. 예상 소요 6분.}
\end{frame}

\begin{frame}{2.3 3분류의 기준 — 58사에게도 가치가 있는가, 커널과 충돌하는가}
\fig[0.56]{fig_2_2_requests.pdf}
\figcap{그림 2-2. 제품화 14건 27 mm / 유상 19건 59 mm / 거절 14건 52 mm. 거절 52 중 ②가 25.}
\keymsg{거절 14건 = 커널 충돌 3 + 마스터 부재 4 + 온담·P1 소관 4 + 축소·별지 3 — C-30(1 mm)이 커널을 깬다}
\note{교재 2.3절. 제품화 14건은 자동차에도 있는 것(처분 코드·건 단위 IQC·규격 대조·계측기·오프라인·SSO·감사 로그·계정 관리·RCA 증빙). 분류가 둘이면(거절 없음) 유상이 거절의 완곡어가 된다. 예상 소요 5분.}
\end{frame}

\begin{frame}{2.4 용량 잠식 — 돈을 받아도 시간은 돌아오지 않는다}
\fig[0.56]{fig_2_3_capacity.pdf}
\figcap{그림 2-3. 168 mm 분모. 전부 138 = 82.1\% / 분류 후 86 = 51.2\% / 유상만 59 = 35.1\% / 거절 0. 상한 6 FTE 선과 에필로그 98 mm 선(봉인).}
\keymsg{분류 후에도 절반 — 상한 6 FTE(연 72 mm)는 "얼마를 받느냐"가 아니라 "얼마를 잃느냐"의 조항}
\note{교재 2.4절 "작은 숫자로 먼저" — 14 × 12 = 168; 138 ÷ 14 = 9.9; 138 ÷ 168 = 82.1. 조현우의 문장(D+230)이 회사 규모로 실현되는 시나리오. ㊱ 항목 5 마지막 줄 "에필로그 예측"은 학생이 봉인 전에 적는다 — 이 그림의 98 mm 선은 학생용 슬라이드에서 가려도 된다. 예상 소요 4분.}
\end{frame}

\begin{frame}{2.5 분류별 계약 문구 — 시점 미약속 · 재사용권 · 누가 하는지}
\begin{tbl}[\footnotesize]{L{22mm}L{72mm}L{36mm}}
분류 & 계약 문구의 핵심 & 리스크 \\ \midrule
제품화 14 · 27 mm & 로드맵 편입 · 무상 · \textbf{릴리스 시점을 약속하지 않는다} & 온담이 Now로 기억(태산 62항목) \\
유상 19 · 59 mm & 별도 SOW 59 × 1,200만 = \textbf{7.08억} · 코드 온담 / 모델 딥게이지 / \textbf{재사용권 딥게이지} · 유지보수 18\% & 온담은 47건을 "21억 포함"으로 읽고 있다 \\
거절 14 · 52 mm & 마스터·법정 DB·IoT·P1·DR은 \textbf{온담 소관} · 무인·감시는 원칙 충돌 · 이물 판정은 측정 후 & "리소스 부족"이면 "돈을 더 주면"이 된다 \\
\end{tbl}
\keymsg{재사용권 한 줄이 없으면 커스텀은 다른 식품 고객에게 팔 수 없는 순수 비용이다}
\note{교재 2.5절. SOW 없이 착수하면 138 mm 전부가 무상. 모델 가중치 소유(온담 데이터로 학습한 모델)는 ㊲ 항목 3 학습 데이터 행. 예상 소요 4분.}
\end{frame}

\begin{frame}{2.6 즉시 탈락 셋 — 47건 목록 위에서는 보이지 않는다}
\begin{enumerate}
\item \textbf{5차원 표 없이 47건 개별 판정} — 판정이 저자를 따라가고, 48번째 요구에 기준이 없다
\item \textbf{②(마스터 없음)를 제품화·유상 커스텀으로} — 판정 근거의 소유가 뒤집히고, 커버리지 38.6\%의 비용이 우리 것이 되며, 3년간 우리가 고친다
\item \textbf{④(법정 기준)를 ㊲ 배상 메모와 연동하지 않음} — 플래그 하나(C-03, 3 mm)가 규제 책임을 옮기고 그 위에 배상 100\%가 얹힌다
\end{enumerate}
\keymsg{C-05는 목록에서 7 mm짜리 한 줄이고 표 위에서는 38.6\%다 — 그래서 표가 먼저다}
\note{교재 2.6절 · 워크북 §1.6 판정 규칙. 즉시 탈락은 방어 불가능한 답이고 그 밖은 근거가 있으면 통과 — 경계 항목(C-19·22·40)의 분류, 상한 FTE 4\til{}7, Go/No-go 결론. 예상 소요 3분.}
\end{frame}

\begin{frame}{1교시 정리 — 제1·2장 한 줄씩}
\begin{itemize}
\item 두 시점 · 21억 = 3.00(6.12) + 11.34(3.78 = ARR) + 6.66(통과) · 라인 70 중 30 정책 밖
\item 네 비율 4.906 / 37.0 / 24.1 / 82.1 · BATNA 대한기공 2.4 × 3 병행 체결 · 답 "벗어나지 않는다"
\item 5차원 표 먼저(40분) · 3분류 27 / 59 / 52 · 용량 168 분모 → 상한 6 FTE · 문구 셋 · 즉시 탈락 셋
\end{itemize}
\keymsg{2교시 — ㊱ 5차원 표의 1차 판정 열을 여러분이 쓴다}
\note{교재 제1·2장 핵심 정리. 예상 소요 2분.}
\end{frame}

% =====================================================================
\section{2교시 (90분) — 실습 ①-a: ㊱ 엔터프라이즈 딜 데스크 시트 · 경계 3영역 · SLA 3층 · AI 판정 오류}
% =====================================================================
\begin{frame}{이 교시에서 하는 것}
\begin{itemize}
\item \textbf{㊱ 딜 데스크 시트} — 워크북 §1.4에 항목 2(5차원 표 — \textbf{1차 판정 열}) · 3(47건 중 표본 12행) · 5(용량 잠식 — 분류 후까지 + \textbf{에필로그 예측 한 줄}) — \textbf{45분}
\item 이어서 제3장 3.1\til{}3.3절 — 경계 3영역 · SLA 3층을 공급자 쪽에서 · AI 판정 오류를 SLA에 넣을 것인가 — \textbf{40분}
\item 순서 — 읽는다 → 찾는다(◆ 카드: 47건 별지 §1.5, 단가 1,200만/mm, 온담의 "21억 포함" 공문) → 쓴다 → 대조한다(X.5·X.6은 다 쓴 뒤)
\end{itemize}
\keymsg{정답을 보지 말고 쓰고, 다 쓴 뒤 대조하라 — 항목 2의 다섯째 열이 비어 있으면 별지는 읽을 수 없다}
\note{워크북 §0.1 · §1.3 카드. 45 + 40 + 정리 5. 표본 12행은 나영선 5 · 한동석 3 · 박세연 2 · 오정근 1 · 최영주 1. 예상 소요 2분.}
\end{frame}

\begin{frame}{㊱ 딜 데스크 시트 — 과제 지시}
\begin{itemize}
\item \textbf{과제}: 항목 2(차원 다섯 × 1차 판정·근거 — 근거는 "58사 가치"와 "커널 충돌" 둘뿐) · 3(표본 12행 — 차원 태그를 옮겨 적고 예외는 근거 열에) · 5(4 시나리오 — mm ÷ 14 · ÷ 168 · 릴리스 공백 · 판단 + 에필로그 예측)
\item 시점 \textbf{D+800}. 카드가 주는 것 — 정본 §5.2\til{}5.8 · E-01\til{}17 · G-01 개발 14 · 47건 별지(차원·공수·부담은 [추가 설정]) · 단가 1,200만/mm · 원가 1,000만/mm
\item 앞 산출물 — ㉖ 항목 5 Won't(B-17·18·20) · ㉟ 항목 4·5(CR-044 6.0 mm) · ㉕ 커널 6-1·6-2 · ㉓ 초기/확장 SAM · ㉘ 라인 = 밸류메트릭
\item 여러분이 결정하는 것 — 1차 판정 열 · 분류 · 상한 FTE · Go/No-go의 조건 · "벗어날 것인가"의 답
\end{itemize}
\keymsg{항목 5 마지막 줄 "에필로그 예측"을 지금 적는다 — 마무리 30분에 대조한다}
\note{워크북 §1.1\til{}1.3 + ◆ 카드. 온담 공문의 "21억 포함" 문구를 카드에 넣은 이유 — 유상 커스텀 SOW의 수용 여부가 미정임을 학생이 알아야 항목 4의 리스크 열을 쓴다. 예상 소요 5분.}
\end{frame}

\begin{frame}{작성 요령 ㊱ — 차원 · 별지 · 용량 (항목 2·3·5)}
\begin{itemize}
\item \textbf{항목 2}: 차원 단위로 셋 중 하나이되 한 차원 안에서 갈릴 수 있다 — 갈리면 C-ID로. ②는 거절, ④는 "㊲ 항목 5 연동"을 문장에
\item \textbf{항목 3}: 차원이 같으면 같은 분류가 기본 · 5차원 밖은 엔터프라이즈 공통(제품화)인가 온담 전용(유상)인가 · 저자별 합을 마지막에 검산(19/11/9/5/3 · 62/31/28/12/5)
\item \textbf{항목 5}: 분모 168을 먼저 · 제품화도 공수다(값을 안 받을 뿐) · 거절 행에 대한기공 2.4를 적는다 · 예측 한 줄
\item 유지보수 부담(H·M·L)을 비우지 않는다 — 커스텀은 만드는 날이 아니라 3년간 고치는 날이 비용
\end{itemize}
\keymsg{탈락 조건 — 표 없이 개별 판정 · ②를 우리가 · ④ 미연동 · 분류가 둘 · 21억 또는 7억을 ARR로}
\note{워크북 §1.3 가이드 + §1.6 판정 규칙. 정답 처리 — 경계 항목의 다른 분류(근거 열 있으면) · 상한 4\til{}7 FTE · Go/No-go 결론(조건이 계약 문장이면). 예상 소요 5분.}
\end{frame}

\begin{frame}{㊱ 흔한 오류 점검}
\begin{itemize}
\item 항목 2를 건너뛰고 별지부터 채웠다 → 저자별로 판정이 흔들린다 → 탈락
\item C-05·06·07·36을 "고객이 원하니 유상"으로 → 마스터의 소유가 뒤집힌다 → 탈락
\item C-03을 유상으로 받고 ㊲ 항목 5를 언급하지 않았다 → 3 mm가 7억이 된다 → 탈락
\item 138 mm를 "1년이면 된다"로 · 제품화를 공수에서 뺐다 · 거절 행에 매출 0만 적고 2.4를 안 적었다
\item Go 조건이 "신중하게" — 계약 문장이 아니면 조건이 아니다
\end{itemize}
\keymsg{C-30 주말·야간 무인 판정은 1 mm다 — 작아서 받으면 커널이 깨진다}
\note{워크북 §1.6 해설. 조별 순회 시 항목 2의 다섯째 열이 채워졌는지, ②가 거절인지를 먼저 본다. 예상 소요 3분.}
\end{frame}

\begin{frame}{예시본 참조 — 워크북 §1.5\til{}1.6}
\begin{itemize}
\item 항목 1 딜 개요 7행(정본 ID) · 항목 2 5차원 + 집계(①24 ②25 ③26 ④8 ⑤12 —43) · 항목 3 표본 12 + 별지 47 · 항목 4 문구 셋 · 항목 5 138 / 86 / 59 / 0 · 항목 6 3표기 · 항목 7 BATNA 7행 · 항목 8 조건부 Go 일곱 조건 + 대한기공 병행
\item 조건부 Go — ① 3분류 별지 ② 상한 6 FTE ③ SOW 7.08억 ④ 정본 판정자 나영선 ⑤ 측정 80건 ⑥ 배상 4구간 ⑦ 라인 별지. ①②⑤ 중 하나라도 빠지면 No-go
\item 대조 규칙 — 항목 2는 ②·④의 판정, 항목 3은 저자별 합, 항목 5는 168 분모의 유무
\end{itemize}
\keymsg{인쇄용 완성 예시 — 00\_Day8\_템플릿팩\_완성예시.pdf(16쪽)}
\note{예시본은 저녁에. 수업 중에는 판정 규칙 상자만. 예상 소요 3분.}
\end{frame}

\begin{frame}{3.1 경계 3영역 — 거절은 우리 언어, 합의서는 누가 하는지}
\fig[0.54]{fig_3_1_boundary.pdf}
\figcap{그림 3-1. 제품 14 / 유상 19 / 자체 8 + 미확정 5 → 3자 회의에서 0. 학습 데이터 행 — 측정 세트 80건의 라벨 판정자 = 정본 판정자.}
\keymsg{미확정을 남기고 서명하지 않는다 · 학습 데이터 행이 없으면 ㊵의 측정 세트가 누구 것인지 정해지지 않는다}
\note{교재 3.1절 · ㊲ 항목 2·3. 유상의 소유권 세 줄(코드 온담 / 모델 딥게이지 / 재사용권 딥게이지)과 유지보수 18\%. 예상 소요 5분.}
\end{frame}

\begin{frame}{3.2 SLA/OLA/UC — 공급자 쪽에서. 온프렘은 고객이 UC 당사자}
\fig[0.52]{fig_3_2_sla_chain.pdf}
\figcap{그림 3-2. 여섯 행 중 둘이 문제 — P1 야간 4.5h > 4h(불성립), AI 판정 정확도(모집단 없음). 온담 IT본부가 UC인 층은 진하게.}
\keymsg{Day4 ⑰에서 수행사의 4시간 UC를 거부한 사람이 오늘 그 수행사다 — 야간은 강등하거나 온담이 OLA를 서명한다}
\note{교재 3.2절 · E-13·14. 219.2분 = 30.4375 × 24 × 60 × 0.5\%. GPU 2대는 성능이 아니라 가용성 조건. 릴리스 행 "온담 승인 지연 기간의 판정 오류는 SLA 제외". 예상 소요 6분.}
\end{frame}

\begin{frame}{3.3 AI 판정 오류를 SLA에 넣을 것인가 — 넣되 한정한다}
\begin{itemize}
\item 온담은 넣자고 한다 — 정당하다, 그것이 사는 물건이니까
\item 문제 ① \textbf{모집단}: FN ≤ 3.0\%는 자동차 3세트 60건의 기준 — 온담 세트에서 측정된 적이 없다. 모집단 없는 SLA 항목은 무효
\item 문제 ② \textbf{연동}: SLA에 넣는 순간 배상 조항이 잡는다 — 61점 중 한 점포의 오판정 하나가 연 청구 7억의 배상 사유
\item 답: 측정 세트에서 측정된 FN에 한정 · 위반 시 서비스 크레딧 · 배상 대상 아님 · 자동판정 기본 OFF · \textbf{최종 판정은 온담 검사원}
\end{itemize}
\keymsg{세영 사고의 로그가 "최종 판정은 검사원"을 입증했다 — 커널 6-1은 원칙이면서 배상 조항이다}
\note{교재 3.3절 · V-06·07 · Day7 ㉞ 항목 4. SLA 표 여섯째 행의 문구 — "온담 4세트 미측정 — 세트가 없으면 이 행 자체가 성립하지 않는다 / 측정 계획(㊵ 항목 4) 완료 전 임시 기준". 예상 소요 5분.}
\end{frame}

\begin{frame}{2교시 정리 — ㊱과 제3장 앞부분}
\begin{itemize}
\item \textbf{차원 표가 본체} — 별지는 사본. ②는 우리가 만들지 않는다, ④는 배상과 연동한다
\item \textbf{상한 6 FTE} — 돈이 아니라 시간의 조항. 에필로그 예측을 적어 두었는가
\item \textbf{경계 3영역} — 거절 → 자체 개발(누가 하는지). 미확정 0 · 학습 데이터 행
\item \textbf{SLA 3층 공급자 판} — 야간 P1 불성립, 온담이 UC 당사자, AI 판정은 측정 세트 한정·크레딧·배상 제외
\end{itemize}
\keymsg{3교시 — 배상의 물리적 하한(6.12로), 218문항, 이관의 정본 판정자, 그리고 ㊲ 실습}
\note{2교시 마무리. 예상 소요 2분.}
\end{frame}

% =====================================================================
\section{3교시 (90분) — 배상 4구간 · 218문항 · 데이터 이관 · 실습 ①-b: ㊲ 경계 합의서·SLA·배상 메모}
% =====================================================================
\begin{frame}{이 교시에서 배우는 것}
\begin{itemize}
\item \textbf{배상 상한} — 물리적 하한은 검산값(6.12)으로: 1년차 $-$4.04 · 3년 공헌 5.52 · 100\% = 7.0 = 공헌의 127\% · 4구간은 조건으로
\item \textbf{218문항} — 받는 쪽의 SAC: 즉시 131 / 조건부 62 / 불가 25 · 제218문항은 ㉞ 전문
\item \textbf{이관} — 정본 판정자는 사람(나영선) · 커버리지 61.4\% 1차 → 마스터 후 95\% · 38.6\%는 온담 것 · 3.00은 61.4\% 기준
\item \textbf{실습 ①-b} ㊲ 항목 4(P1·AI 판정 2행) · 5(물리적 하한 + 4구간) · 7(정본 판정자 + 38.6\%) — 35분
\end{itemize}
\keymsg{"50\%에서 만나자"는 조건 없는 중간이다 — 다음 실사에서 "왜 50\%였나"에 답이 없다}
\note{교재 3.4\til{}3.6절 + ㊲ 실습. 45 + 35 + 정리 10. 예상 소요 2분.}
\end{frame}

\begin{frame}{3.4 배상 상한 — 물리적 하한과 4구간}
\fig[0.55]{fig_3_3_liability.pdf}
\figcap{그림 3-3. 10 / 25 / 50 / 100\% = 0.7 / 1.75 / 3.5 / 7.0억 vs 3년 공헌 5.52 · 현금 7.52. AI 판정 오류·법정 판정은 어느 구간에도 없다.}
\keymsg{100\% = 3년 공헌의 127\% = 현금의 93\% — 한 번이면 회사가 없다. 3.00으로 계산하면 7억이 "감당 가능해 보인다"(F-14)}
\note{교재 3.4절 "작은 숫자로 먼저" — 7.00 $-$ 6.12 $-$ 2.22 $-$ 2.70 = $-$4.04 · 21.00 $-$ 6.12 $-$ 6.66 $-$ 2.70 = 5.52. 4구간 조건 — 일반 SLA(크레딧) / 이관·output 미달 / 보안 사고 / 고의·중과실. 권고: 25\%를 열고 100\%의 조건을 지킨다 — 온담이 두려워하는 것은 이관 실패(11년치)이지 오판정이 아니다. 예상 소요 6분.}
\end{frame}

\begin{frame}{3.5 218문항 — 받는 쪽의 SAC}
\fig[0.5]{fig_3_4_q218.pdf}
\figcap{그림 3-4. 즉시 131(제218문항 포함) / 조건부 62(모의해킹 1,800만·4주) / 불가 25(ISMS-P 2029-Q4 계획 · 인프라·DR은 온담 통제).}
\keymsg{불가를 조건부로 쓰지 않는다 — ISMS-P는 4주에 안 된다. 제218문항에는 요약이 아니라 전문 — 숨긴 것이 실사에서 나온다}
\note{교재 3.5절 · E-12 · §4.10 원문 · Day7 ㉟ 규칙 4. 불가 25의 절반은 온프렘이 답 — 온담이 요구한 배포 형태가 온담의 답. PM ⑯ 항목 8(SAC 12항목·판정자·증빙)의 답변자 판. 예상 소요 4분.}
\end{frame}

\begin{frame}{3.6 데이터 이관 — 정본 판정자와 61.4\%}
\begin{tbl}[\footnotesize]{L{30mm}L{70mm}L{30mm}}
항목 & 문구 & 리스크 \\ \midrule
정본 판정자 & \textbf{나영선}(이천·안성은 위임 판정자, 상충 시 나영선 최종) & 정본은 미지정 — 법인은 판정하지 않는다 \\
검증 6종·통과 & 건수 100 / 합계 0 / 키 100 / 참조 99.5 / \textbf{커버리지 61.4 → 마스터 후 95} / 표본 300건 0.5\% & 95\%를 1차 기준으로 쓰면 이관 실패 조항 \\
38.6\% 미매핑 & 임시 식별자 이관 · 판정 초안 제외 · 온담 마스터 정비 기한 6개월 · 재매핑 별도 & 기한 없으면 영원히 임시 \\
구축비 범위 & 3.00 = 61.4\% 기준(11,200권 · 4,100개) · 범위 밖 별도 & $-$4.04가 $-$7.16이 된다 \\
\end{tbl}
\keymsg{②의 즉시 탈락이 여기서 재발한다 — 38.6\%를 우리가 메우면 구축비는 6.12에서 더 오른다}
\note{교재 3.6절 · 정본 §5.5 · E-03·17. 판정자 후보 셋 중 근거 있는 선택은 통과(조성태 — 별도 법인의 권리). 즉시 탈락은 "온담"이라 적는 것. 예상 소요 4분.}
\end{frame}

\begin{frame}{㊲ 경계 합의서·SLA·배상 메모 — 과제 지시}
\begin{itemize}
\item \textbf{과제}: 항목 4(SLA 6행 중 P1 복구·AI 판정 정확도 — OLA / UC / 사슬 합 / 성립?) · 5(물리적 하한 계산 → 4구간 조건 → 권고) · 7(정본 판정자 · 38.6\% 처리) — \textbf{35분}
\item 시점 \textbf{D+814} 3자 회의. 카드 — 0.7 FTE 평일 주간 · 야간 대기 0.5 · NBD · 온담 야간 1차 30분 · 218 회신 131/62/25 · 1년차 $-$4.04 · 3년 공헌 5.52 · 현금 7.52
\item 앞 산출물 — ㊱ 항목 3·4 · PM ⑰ 항목 3(219.2분)·4(3층 규칙) · PM ⑯ 검증 6종 · ㉞ 항목 4(162만)·8(증빙)·7(미감시) · ㉘ 항목 6 · ㉗ 항목 1
\item 여러분이 결정하는 것 — 미확정 5건 · SLA 성립 판정과 문구 · AI 판정의 처리 · 4구간 조건 · 정본 판정자
\end{itemize}
\keymsg{탈락 조건 — 하한을 3.00으로 · AI 판정을 배상 구간에 · 법정 판정 미분리 · 판정자를 법인명으로 · 미확정 남기고 서명 · 분(219.2) 없음}
\note{워크북 §2.3 카드 + §2.6 판정 규칙. 정답 처리 — P1 야간의 처리 셋(강등 / 온담 OLA / 대기 증원 + 비용) · 4구간 경계가 조건과 함께면 · 판정자 후보 셋 중 근거 있는 선택. 예상 소요 5분.}
\end{frame}

\begin{frame}{작성 요령 ㊲ — 사슬 · 하한 · 판정자 (항목 4·5·7)}
\begin{itemize}
\item \textbf{항목 4}: 합산은 가장 긴 사슬 · 한 층이 비면 불성립 · 고객이 서명해야 하는 층을 우리 OLA로 쓰지 않는다 · AI 판정 행은 "세트가 없으면 행이 없다"를 적고 ㊵ 항목 4로 보낸다
\item \textbf{항목 5}: 6.12로 계산 → 1년차 $-$4.04 · 3년 공헌 5.52 · 현금 7.52 → 100\% = 127\% / 93\% → 4구간을 조건으로 → AI 판정·법정 판정은 어느 구간에도 없다 → 권고(25\%를 열고 100\% 조건을 지킨다)
\item \textbf{항목 7}: 판정자 한 사람 · 6종 통과 기준 수치 · 38.6\% 임시 식별자 + 온담 정비 기한 · 구축비 범위 61.4\% 기준
\item 항목 2의 미확정 5건은 §2.5로 확인하되 "서명 전 0"을 항목 8에 적는다
\end{itemize}
\keymsg{같은 조항을 반대편에서 — ⑰의 배상 조항이 오늘 우리를 향한다. 선례 162만은 "상한은 지키고 서비스로 관계를 산다"}
\note{워크북 §2.3 항목 4·5·7 가이드. 예상 소요 5분.}
\end{frame}

\begin{frame}{㊲ 흔한 오류 점검 · 예시본 참조}
\begin{itemize}
\item 야간 P1을 0.7 FTE로 성립시켰다(⑰에서 여러분이 거부한 것) · 99.5\%만 있고 분이 없다
\item AI 판정 정확도를 안 넣었다(온담이 계약 안 한다) 또는 모집단 없이 넣었다 또는 배상에 연동했다
\item "50\%에서 만나자" · 3.00으로 하한 · C-11을 받은 채 법정 판정을 우리가
\item 판정자를 "온담"으로 · 95\%를 1차 통과 기준으로 · 마스터 정비 기한 없음
\item 예시본 §2.5 — 당사자 3 · 3영역 + 미확정 0 · 소유권 4행 · SLA 6행 · 배상 6행(4구간) · 218 3행 · 이관 5행 · 서명 3 + 입회 2
\end{itemize}
\keymsg{인쇄용 완성 예시 — 37\_경계합의서SLA배상\_완성예시.pdf(3쪽)}
\note{워크북 §2.6 해설. 조별 순회 시 항목 5의 하한 계산에 6.12가 있는지 먼저 본다. 예상 소요 3분.}
\end{frame}

\begin{frame}{3교시 정리 — 제3장 한 줄씩}
\begin{itemize}
\item 경계 3영역 — 제품 14 / 유상 19 / 자체 8 · 미확정 0 · 학습 데이터 행
\item SLA 3층 — 공급자 쪽 · 야간 P1 불성립 · 온담이 UC · AI 판정은 측정 세트 한정
\item 배상 — 6.12로 하한(F-14) · $-$4.04 / 5.52 / 7.52 · 4구간은 조건 · 25\%를 열고 100\%의 조건을 지킨다
\item 218문항 — 131 / 62 / 25 · 제218문항 전문 · 이관 — 판정자는 사람 · 61.4 → 95 · 38.6\%는 온담 것
\end{itemize}
\keymsg{오후 — 자본(18.9 vs 15.68 · 1.29\%p · 2×2)과 편익 원장(output 열 · 모집단)}
\note{3교시 마무리. 예상 소요 2분.}
\end{frame}

% =====================================================================
\section{4교시 (90분) — 자본: IR 두 값 · 3표기 · 3케이스 · 캡테이블 · 옵션풀 1.29\%p · 레드팀·2×2 · 실습 ② ㊳·㊴}
% =====================================================================
\begin{frame}{이 교시에서 배우는 것}
\begin{itemize}
\item \textbf{IR 1페이저} — 검산값을 앞에, 계약 기준을 괄호에, 분모를 각주에 · 커널 한 줄이 채택률의 답 · 덱은 ㉝ 규칙을 절반만 적용했다
\item \textbf{3케이스} — 결정으로 가른다(수용 / 거절 / 확장) · 80억은 6개월차 · Bear = 코너스톤 소멸 + pre 242 · \textbf{IC 8문} — 답에 분모
\item \textbf{캡테이블} — 주식수로. T2 1,305,556 → A 1,772,879 / B 1,665,706 · \textbf{옵션풀} 같은 112,745주, pre 20.00 / post 18.71 → 1.29\%p = 25.4억
\item \textbf{레드팀·2×2} — 참가적 40.6억 · drag 34.1\% · 정보권 ARR 정의 · 거절 × 코너스톤 소멸 · 거절 × 한강 pre 242
\item \textbf{실습 ②} ㊳ 항목 2·6(25분) · ㊴ 항목 3·6(25분)
\end{itemize}
\keymsg{"표준 문구예요"는 참이다 — 표준이라는 것과 값이 없다는 것은 다르다}
\note{교재 제4장 + ㊳·㊴ 실습. 이론 40 + 25 + 25. 예상 소요 2분.}
\end{frame}

\begin{frame}{4.1\til{}4.2 IR 1페이저 — 두 값 병기 · 집중 리스크 3표기}
\fig[0.55]{fig_4_1_arr_three.pdf}
\figcap{그림 4-1. ARR 18.9 / 15.68 / 11.9 — 18.9에는 구축 3.0 + 운영 2.22가 들어 있다. 집중 리스크 37.0(과대) / 24.1(적정) / 82.1(덱에 없다).}
\keymsg{검산을 앞에 쓰면 "자기 숫자를 깎아서 내는 회사"가 되고, 그것이 pre 320(검산 20.4×)의 근거다 — IC가 15.68로 다시 계산하기 전에}
\note{교재 4.1·4.2절 · ㊳ 항목 1·2. 덱의 절반 적용 — GM은 계약 MRR(54), burn multiple은 순SaaS(1.20): 각각 더 좋아 보이는 분모. 서영채 D+598 메일이 D+802 초안 검토에서 재현되기 전에. 예상 소요 5분.}
\end{frame}

\begin{frame}{4.3 재무모델 3케이스 — 결정으로 가른다}
\fig[0.56]{fig_4_2_three_cases.pdf}
\figcap{그림 4-2. Bear(온담 거절 + 대한기공) 16.8 → 19.2 / Base(3분류 수용) 21.8 → 30.3 / Bull(+ 식품 3사) 25.9 → 43.3. 현금은 6개월차 2.4억이 바닥.}
\keymsg{거절이 Bear인 이유는 매출이 아니라 코너스톤 소멸 + 한강 재협상 — "거절하면 죽는다"가 아니라 "거절하면 시리즈B를 앞당긴다"}
\note{교재 4.3절 · ㊳ 항목 3·4. 클로징 전 6개월 0.85억 × 6 = 5.1 → 2.4억. 사용처 80 = 32 + 20 + 8 + 5.82(온담 적자 3.12 + 제품화 2.70 — 숨기지 않는다) + 14.18. 예상 소요 5분.}
\end{frame}

\begin{frame}{4.4 IC 반대신문 8문 — 답에 분모를}
\begin{tbl}[\scriptsize]{L{48mm}L{82mm}}
질문 & 답(분모 포함) \\ \midrule
집중 리스크가 37\%인데 24.1\%라고요? & 18.9는 청구 전액, 15.68은 라이선스만. 그리고 \textbf{82.1\%}(138/168)를 우리가 먼저 \\
채택률 26\%면 안 쓰이는 것 아닌가요? & 분모 ㉮. 기본 OFF 설계 — 커널. 지표는 NS-1(58사 27,600건/주). FP 21.3이 원인, DS-v2가 조치 \\
온담을 거절하면 코너스톤은? & 소멸. 한강 단독 → 풀 협상력 감소. BATNA 대한기공 2.4 × 3 \\
GM 54\%인데 원가율 44\%면? & 9,917 vs 8,583. 같은 분모면 46.9 / 38.1. 병기했다 \\
사고 이후 재발했습니까? & 0건 — 감시 대시보드 위의 0건. 재발 방지 6건 증빙 = 제218문항 답변 \\
구축비 3억이 6.12억이면 적자 아닌가요? & 1년차 $-$4.04 · 3년 공헌 5.52(SOW 7.08 별도) → 상한 10\%·커스텀 유상 \\
\end{tbl}
\keymsg{답이 "분모가 다르다"로 끝나면 안 된다 — 어느 것이 적정인지, 그리고 IC가 아직 묻지 않은 숫자}
\note{교재 4.4절 · ㊳ 항목 6(8문 중 6). 나머지 둘 — 온담 투입 인원(86 mm ÷ 12 = 7.2 FTE → 상한 6) · NDR 104의 분모(D+435 21사 2,150). "검사원 교육"은 ㉞의 즉시 탈락. 예상 소요 4분.}
\end{frame}

\begin{frame}{4.5 캡테이블 T2 \ra{} T3 — 주식수로 계산한다}
\fig[0.55]{fig_4_3_captable.pdf}
\figcap{그림 4-3. T2 = 1,250,000 + SAFE 55,556(8억 ÷ 14,400) = 1,305,556. A안 풀 +112,745 → 주당 22,562 → 한강 354,578 → 1,772,879. B안 주당 22,213 → 코너스톤 360,150 → 1,665,706.}
\keymsg{같은 80억 — A안 20.00\% / B안 21.62\%. 강민수 25.38 / 27.01. 차이 1.63\%p는 pre 30억과 풀 112,745주가 합쳐진 것 — 떼어 보는 것이 다음 장}
\note{교재 4.5절 "작은 숫자로 먼저" · K-03\til{}05. 풀 15\% pre: (100,000 + x) ÷ (1,305,556 + x) = 0.15 → x = 112,745. SAFE는 cap 전환(18,050 > 14,400). 심화 — 정본은 cap ÷ FD 1,250,000(pre 방식), 문서가 post-money cap이면 4.44\%. 예상 소요 5분.}
\end{frame}

\begin{frame}{4.6 옵션풀 pre/post — 1.29\%p}
\fig[0.56]{fig_4_4_pool_position.pdf}
\figcap{그림 4-4. 같은 112,745주를 post에 두면 주당 24,511 · 한강 326,384주 · 18.71\%. pre면 22,562 · 354,578주 · 20.00\%. 차이 28,194주 = 1.29\%p = 25.4억.}
\keymsg{"풀 pre"는 pre 320을 294.6으로 깎은 것과 같다 — 요청 문구: "post-money 기준" 또는 "pre 345.4억"}
\note{교재 4.6절 · ㊴ 항목 3. post 방식의 정의(320억 pre를 풀 없는 FD에 적용, 한강도 희석 분담)는 정본 §7.5와 같다. 정의가 적혀 있으면 다른 정의로 다른 숫자도 통과. 시드 0.43\%p → 1.29\%p: 풀 4배 · 밸류 6.7배(Day6 점검 질문 2의 답). 예상 소요 5분.}
\end{frame}

\begin{frame}{4.7 텀시트 레드팀과 2×2}
\fig[0.5]{fig_4_5_two_by_two.pdf}
\figcap{그림 4-5. 거절 × 코너스톤 = 소멸(정본). 거절 × 한강 = 가능하나 IR ARR 18.9 → 14.3, 같은 배수면 pre 242(계산해야 보인다).}
\keymsg{레드팀 — 참가적: 매각 500억 시 창업자 3인 207.1 vs 247.7억 · drag: 한강 + 노들 34.1\%가 우선주 과반 · 정보권: "계약 ARR" 보고가 서영채 D+330을 재현}
\note{교재 4.7절 · ㊴ 항목 1·4·6. 두 텀시트의 기타 조항(참가·drag·이사회·클로징)은 [추가 설정]. 14.3을 만드는 2.4가 대한기공 — ㊱ 조건부 Go가 병행 체결을 포함하는 이유. 예상 소요 4분.}
\end{frame}

\begin{frame}{4.8 채용 여력과 협상 요청 — 셋이면 협상, 다섯이면 거절}
\begin{itemize}
\item 가용 풀 A 172,745주(9.74\%) / B 60,000주(3.60\%) · 27 → 60명 34명 필요 5.8\% → A +69,918주(3.9\%p 여유) / B $-$36,611주(2.2\%p 부족)
\item B의 부족은 시리즈B 전 풀 증설 — 그 희석은 창업자와 코너스톤이 나눈다. 코너스톤이 풀 조항을 안 넣은 것은 관대함이 아닐 수 있다
\item 요청 ① 풀 post 또는 pre 345.4억 ② 1x 비참가(시드 RCPS 동반 전환, +40.6억) ③ (B안 시) 선행 조건을 "ARR 3.78억 이상 신규 계약"으로
\item BATNA 정직하게 — 다른 리드 없음 / 브릿지(서영채 자기 희석 반대) / 매출로 버티기 불가(8.9개월)
\end{itemize}
\keymsg{권고: 한강 A안 + 요청 1·2 — 온담 조건부 수용·대한기공 병행. 코너스톤은 요청 3이 관철될 때만}
\note{교재 4.8절 · ㊴ 항목 5·7. 직급별 옵션 0.40 / 0.15 / 0.05 / 1.00(첫 PM 0.4\% 기준). 두 텀시트는 같은 문제(풀)를 반대 시점에 푼다. 예상 소요 4분.}
\end{frame}

\begin{frame}{실습 ② — ㊳ 항목 2·6 (25분) · ㊴ 항목 3·6 (25분) 과제 지시}
\begin{itemize}
\item \textbf{㊳}: 항목 2(트랙션 표 — ARR·집중 리스크·payback 3행: IR 기재 / 검산 / 분모·각주) · 6(IC 8문 중 1·3·6 — 집중 리스크·코너스톤·구축비 적자)
\item \textbf{㊴}: 항목 3(옵션풀 pre/post — \textbf{주식수로} 1.29\%p를 만든다) · 6(2×2 — 거절 × 한강 칸의 pre 재산정 포함)
\item 카드(㊳) — COGS 4,560 · NDR 기준 D+435 21사 2,150 · CAC 740 · 케이스별 순소진 · 온담 포함/제외 정의. 카드(㊴) — K-02\til{}05 · 기타 조항 · 기부여 40,000주 · post 방식의 정의
\item 여러분이 결정하는 것 — 어느 값을 앞에 쓸지(근거) · IC 답의 문장 · 2×2 네 칸의 판정 · 요청 순위
\end{itemize}
\keymsg{탈락 조건 — 18.9만 기재 · 검산 열 공란 · NDR 기준 코호트 없음 / 풀 효과를 \%로 · SAFE 전환 누락 · 거절 × 코너스톤이 소멸이 아님}
\note{워크북 §3.3·§4.3 카드 + §3.6·§4.6 판정 규칙. 정답 처리 — 검산값을 뒤에 써도 근거 있으면 · post 방식의 다른 정의(정의가 적혀 있으면) · 권고가 코너스톤(요청 3 조건이 있으면). 예상 소요 4분.}
\end{frame}

\begin{frame}{㊳·㊴ 흔한 오류 점검 · 예시본 참조}
\begin{itemize}
\item ㊳ — 18.9만 / 두 값을 쓰되 분모 없음 / 온담을 세 케이스 모두에 / 80억을 0개월차에 / 채택률 답에 "검사원 교육" / 온담 적자를 운전자본에 숨김
\item ㊴ — \%로 계산(풀이 어디 있는지 안 보인다) / SAFE를 discount로 전환(cap이 더 낮다) / 2×2를 O/X로 / 요청 다섯 개 / 정보권 ARR 정의 미검토
\item 예시본 §3.5 — 1페이저 6칸 · 트랙션 6행 · 3케이스 7행 · 사용처 80 · 리스크 5 · IC 8문 / §4.5 — 대조 7행 · 캡테이블 8행 · pre/post 5행 · 레드팀 7행 · 채용 4행 · 2×2 · 요청 3 + 권고
\end{itemize}
\keymsg{인쇄용 완성 예시 — 38\_IR1페이저재무모델\_완성예시.pdf(2쪽) · 39\_텀시트레드팀캡테이블v2\_완성예시.pdf(2쪽)}
\note{워크북 §3.6·§4.6 해설. 조별 순회 시 ㊴ 항목 3에 주식수(354,578 / 326,384)가 있는지 먼저 본다. 예상 소요 3분.}
\end{frame}

\begin{frame}{4교시 정리 — 제4장 한 줄씩}
\begin{itemize}
\item 두 값 병기 — 검산 앞, 계약 괄호, 분모 각주 · 3표기 37.0 / 24.1 / 82.1 · 3케이스는 결정으로 · IC 답에 분모
\item 주식수로 — T2 1,305,556 · A 1,772,879(22,562) · B 1,665,706(22,213) · 강민수 25.38 / 27.01
\item 옵션풀 1.29\%p = 25.4억 · 요청 post 또는 pre 345.4 · 참가적 40.6억 · 거절 × 코너스톤 소멸 · 거절 × 한강 pre 242
\item 채용 A +3.9\%p / B $-$2.2\%p · 요청 셋 · BATNA 정직 · 권고 한강 + 1·2
\end{itemize}
\keymsg{5교시 — ⑳의 표에 여러분 이름이 들어간다. 그리고 인수 기준의 모집단}
\note{4교시 마무리. 예상 소요 2분.}
\end{frame}

% =====================================================================
\section{5교시 (75분) — 편익 원장 공동 확정: ⑳를 반대편에서 · 세 서명 · 인수 기준의 모집단 · 두 시계 · 실습 ③ ㊵}
% =====================================================================
\begin{frame}{이 교시에서 배우는 것}
\begin{itemize}
\item \textbf{⑳를 반대편에서} — 표를 바꾸지 않고 output 열만 더한다 · 소유는 재무본부 · 분모 A는 손대지 않는다 · B 목표 신설(⑳가 예약한 한 가지)
\item \textbf{세 서명} — output 딥게이지 / outcome 한동석·오정근·나영선 / utilisation 실명 + 마스터 정비(온담)
\item \textbf{인수 기준의 모집단} — V-07 그대로 적는다 · 4세트 80건 · 세트별 지표 · 8주 · 1,200만 우리 부담 · 측정 전 OFF
\item \textbf{두 시계} · output 미달만 배상 25\% · \textbf{사슬} — 분모 A 3문서 불변, 자동차 60건 5문서 불변, 기준 2회 변경
\item \textbf{실습 ③} ㊵ 항목 2(B-02(B) 행)·4(모집단 다섯 행) — 25분
\end{itemize}
\keymsg{모집단 없는 fit criterion은 무효다 — 항목 4 첫 행이 비면 항목 8은 서명할 수 없다}
\note{교재 제5장 + ㊵ 실습. 45 + 25 + 5. 온담 G5(2028-03-31)는 지났다 — 71\% 경고, 개선 계획의 공급자가 여러분. 예상 소요 2분.}
\end{frame}

\begin{frame}{5.1\til{}5.2 ㊵ = ⑳의 표 + output 열 — 칸마다 실명}
\fig[0.52]{fig_5_1_ledger_ito.pdf}
\figcap{그림 5-1. ⑳ v1.0 6칸 그대로 + output 열. B-02(A) 불변, B-02(B) 84.6 → G5 88.9 → 95.0(M+36) 신설, B-09·10 신설. B-10의 output: "법정 판정 정확도는 output이 아니다".}
\keymsg{⑳의 output 열 "P1 PM·이재현(수행사)"이 오늘 "딥게이지 강민수"다 — 그리고 utilisation 칸에 마스터 정비(C-05, 온담 자체)가 들어간다}
\note{교재 5.1·5.2절 · PM Day4 ⑳ §5.5 항목 1·2·5 · ㊵ 항목 1\til{}3. G5 실적(B-04 2.9 / B-03 2.6 / B-02A 97.1 / B 88.9 · 71\%)·B-09·10·output 열은 [추가 설정]. output에 outcome(재고 정확도 \%)을 쓰면 즉시 탈락. 예상 소요 6분.}
\end{frame}

\begin{frame}{5.3 인수 기준의 모집단 — FN 2.4\%는 어디서 나온 값인가}
\fig[0.56]{fig_5_2_population.pdf}
\figcap{그림 5-2. 자동차 3세트 60건(측정됨, 기준 2회 변경) vs 온담 4세트 80건(계획 — 세트마다 다른 지표). 옮겨지지 않는다.}
\keymsg{4세트에 "FN ≤ 3.0\%"를 그대로 쓰는 것은 모집단만 틀린 것이 아니라 지표도 틀린 것이다 — 냉장 시계열에 OCR은 없다}
\note{교재 5.3절 · V-06·07 · ㊵ 항목 4. D-1 in-range ≥ 99 / D-2 재현율 ≥ 97 / D-3·4 OCR ≥ 90 \& FN ≤ 5 · 측정 전 자동판정 OFF · 확정 4주 내 나영선·오세진 · 1,200만 딥게이지 부담(온담에 청구하면 측정이 미뤄진다). 예상 소요 6분.}
\end{frame}

\begin{frame}{5.4 두 시계와 미달 시 조치 — output 미달만 배상}
\begin{tbl}[\footnotesize]{L{34mm}L{30mm}L{34mm}L{26mm}}
상황 & 책임 주체 & 조치 & 배상 연동 \\ \midrule
output 미달(생성률·커버리지·측정 세트) & 딥게이지 & 재이관·재측정·무상 연장 & \textbf{㊲ 25\% 구간} \\
outcome 미달(output 충족) & 온담 임원 & 개선 계획(⑳ 규칙) & 비연동 \\
utilisation 미달 & 나영선·점포 QA & 교육·체크리스트·점검 & 비연동 \\
측정 불가(분모 B 실사 없음) & 온담 물류본부 & 측정 체계 먼저 — 보류 & 비연동(측정 세트는 output) \\
\end{tbl}
\keymsg{outcome을 배상에 넣으면 ⑳ 항목 5의 책무 분리가 무너지고, output을 배상에서 빼면 ㊲ 25\% 구간의 대상이 없다}
\note{교재 5.4절 · ㊵ 항목 5·6. 네 리뷰 — 원장 v2(2028-05-31, 재무본부 판정, 딥게이지 배석) / 측정 확정(서명 + 8주) / 운영 +24(2029-05) / 분기 릴리스(온담 승인 후). 공급자는 배석이지 판정자가 아니다. 예상 소요 4분.}
\end{frame}

\begin{frame}{㊵ 편익 원장 공동 확정본 — 과제 지시 (25분)}
\begin{itemize}
\item \textbf{과제}: 항목 2(B-02(B) 행 — 측정식·기준값·목표값·시점·책임자·\textbf{output 열}) · 4(모집단 — 현재 기준의 모집단 / 4세트 측정 계획 / 임시 기준·유효 기간 / 확정 절차·판정자 / 비용 부담)
\item 시점 \textbf{D+821} — ⑳ 리뷰 계획의 "2028-05 운영 +12 원장 v2". 카드 — G5 실적 · B-09·10 표본 · 측정 세트 D-1\til{}4 × 20 · 8주 · 1,200만
\item 앞 산출물 — \textbf{⑳ 항목 1·2·4·5 그대로} · ㉗ 항목 1·3 · ㉞ 항목 5·6·7 · ㊲ 항목 5·7 · ㊱ C-27·28
\item 여러분이 결정하는 것 — B 목표값·시점 · output 지표 셋 · 세트별 임시 기준 · 확정 판정자 · 비용 부담
\end{itemize}
\keymsg{탈락 조건 — 항목 4 첫 행 공란 또는 "FN ≤ 3.0\%" 그대로 · 분모 A 목표 변경 · output에 outcome · 법정 판정을 output으로 · outcome 미달을 배상에 · 항목 4 공란인 채 서명}
\note{워크북 §5.3 카드 + §5.6 판정 규칙. 정답 처리 — B 목표 93\til{}96\% · 시점 M+30\til{}36 · 측정 60\til{}120건 · 세트별 기준 수치(근거 있으면). 항목 7(사슬)은 6교시 조별. 예상 소요 4분.}
\end{frame}

\begin{frame}{㊵ 작성 요령 · 흔한 오류 · 예시본 참조}
\begin{itemize}
\item \textbf{항목 2}: ⑳ 문장을 바꾸지 않는다 — 오늘 바꿀 것은 B의 목표 신설뿐. output은 딥게이지가 인도하는 것을 측정 가능한 지표로(자동 생성률 ≥ 90 · 커버리지 61.4 → 95 · 판정 정확도는 미측정 → 항목 4)
\item \textbf{항목 4}: 첫 행은 V-07의 문장 그대로 — 가장 불리한 문장을 먼저. 세트마다 지표가 다르다. 비용은 우리 것
\item 흔한 오류 — 딥게이지 문서로 만듦(소유는 재무본부) · output에 "GAUGE"(제품명은 서명하지 않는다) · 마스터 정비를 output에(② 재발) · 측정 비용을 온담에
\item 예시본 §5.5 — 문서 정보 4행 · 편익 6행 · ITO 5행 · 모집단 5행 · 리뷰 4행 · 미달 4행 · 사슬 6행 · 서명 4
\end{itemize}
\keymsg{인쇄용 완성 예시 — 40\_편익원장공동확정본\_완성예시.pdf(4쪽)}
\note{워크북 §5.3 가이드 · §5.6 해설. 조별 순회 시 항목 4 첫 행에 "온담 4세트 미측정"이 있는지 먼저 본다. 예상 소요 4분.}
\end{frame}

\begin{frame}{5교시 정리 — 제5장 앞부분 한 줄씩}
\begin{itemize}
\item ㊵ = ⑳ + output 열 · 소유 재무본부 · 분모 A 불변 · B 84.6 → 88.9 → 95.0(M+36) · B-09·10 신설
\item 세 서명 — output 딥게이지 / outcome 한동석·오정근·나영선 / utilisation 실명 + 마스터 정비(온담)
\item 모집단 — V-07 그대로 · 80건 · 세트별 지표 · 8주 · 1,200만 · 측정 전 OFF · 확정 나영선·오세진
\item 두 시계 · output 미달만 25\% · 공급자는 배석
\end{itemize}
\keymsg{6교시 — 이 표의 모든 칸이 어디서 왔는가를 거꾸로 걷는다}
\note{5교시 마무리. 예상 소요 2분.}
\end{frame}

% =====================================================================
\section{6교시 (75분) — 실습 ④: 8일 사슬 역추적 (조별) · 발표}
% =====================================================================
\begin{frame}{이 교시에서 하는 것 — ㊵ 항목 7, 40종을 거꾸로}
\begin{itemize}
\item \textbf{조별 40분} — ㊵ 항목 7의 표를 벽에 붙이고 두 사슬을 끝까지 걷는다. 각 칸의 출처 문서·항목·날짜, 그리고 변경 이력(몇 번 바뀌었는가)
\item \textbf{PM 사슬} — 분모 A 91.3: ㊵ 항목 2 ← ⑳ 항목 2(2027-05-21) ← ⑫(2026-09) ← ① BC §4(2025-11)
\item \textbf{Product 사슬} — "자동차 3세트 60건": ㊵ 항목 4 ← ㉞ 항목 5(D+540) ← ㉗ 항목 1(D+137) ← ㉖ 항목 3(D+130) ← ㉓(D+60) ← ㉒ 박선재(D+40)
\item 더 걷는다 — 배상 사슬(㊵ 6 ← ㊲ 5 ← ㉞ 4 ← ㉘ 6 ← ⑰) · 시계 사슬(㊵ 5 ← ⑳ 4 ← ⑱ 이관 6번 ← 2025-11 이사회 반려 · ② 헌장 §5 제외 범위)
\item \textbf{발표 25분} — 조마다 한 문장: 모집단 칸을 비운 채 서명한 문서가 8일 동안 몇 개였는가
\end{itemize}
\keymsg{기준은 여러 번 바뀌었는데 모집단이 한 번도 안 바뀌었다면, 그 기준은 결과에 맞춘 것이다}
\note{워크북 §5.3 항목 7 · §7.2 조별 과제. 힌트 — ㉗ v1.0의 합격 기준, ㉞ v2의 개정, ㊲ 항목 4 초안, IR 덱 18.9. 40 + 25 + 정리 10. 예상 소요 4분.}
\end{frame}

\begin{frame}{5.5 8일 사슬 역추적 — 두 사슬이 만나는 칸}
\fig[0.53]{fig_5_3_chain.pdf}
\figcap{그림 5-3. PM 사슬(분모 A — 세 문서 불변)과 Product 사슬(자동차 60건 — 다섯 문서 불변, 기준 2회 변경)이 ㊵ 항목 4 첫 행에서 만난다.}
\keymsg{"검사기록은 제외 범위"(② 헌장 §5, 2025-11)가 3년 뒤 21억 계약이 됐다 — 첫 문서의 한 줄이 마지막 문서의 계약}
\note{교재 5.5절. 불변인 칸을 "안정적"으로 읽지 않는다 — 모집단 불변은 안정이 아니라 미측정. 조별 순회: 사슬이 한 단계에서 끝난 조는 날짜를 물어 다음 단계로. 예상 소요 4분.}
\end{frame}

\begin{frame}{발표 — 조마다 한 문장}
\begin{itemize}
\item 모집단 칸을 비운 채 서명한 문서 — 몇 개였고, 어느 것이었는가
\item 그중 하나를 골라: 그 칸을 채우는 데 걸렸을 시간은 얼마였는가(측정 세트 80건 = 8주 · 1,200만)
\item 그리고 ㊱ 항목 5에 적어 둔 "에필로그 예측" — 마무리 30분에 연다
\end{itemize}
\keymsg{같은 문장, 8일 전 — 재고 정확도 91.3\%는 순환실사 SKU(69\%)의 값이었다. 그때 감사자, 지금 서명자}
\note{발표 25분(조당 3\til{}4분). 강사는 조의 답을 칠판에 세지 않는다 — 마무리에서 최종 명제와 함께 센다. 예상 소요 2분.}
\end{frame}

\begin{frame}{6교시 정리 — 사슬이 보여 준 두 가지}
\begin{itemize}
\item \textbf{불변인 것} — 분모 A(3문서) · 자동차 60건(5문서). 하나는 정직한 불변(승인된 분모로 판정), 하나는 미측정의 불변
\item \textbf{변한 것} — 합격 기준 두 번(8.0 → 3.0) · 여러분의 자리(감사자 → 서명자) · "제외 범위"가 계약이 된 것
\item 다음 30분 — 봉인을 연다
\end{itemize}
\keymsg{측정 시점에 밝혀진 것은 서명 시점에 감출 수 없다 — 8일 동안 그 문장을 마흔 번 썼다}
\note{6교시 마무리. 예상 소요 2분.}
\end{frame}

% =====================================================================
\section{마무리 (30분) — 에필로그 봉인 해제 · 8일 최종 명제 · 학술 배경 · Product 트랙을 마치며}
% =====================================================================
\begin{frame}{5.6 에필로그 — 2028-11-26 (D+1000), 봉인 해제}
\begin{tbl}[\footnotesize]{L{34mm}L{62mm}L{30mm}}
항목 & 값 & 출처 \\ \midrule
인원 & \textbf{34명}(개발 21 / 세일즈 7 / CS 4 / 경영지원 2) & G-01 · §3.4 \\
유료 / ARR(온담 제외) & 58개사 / \textbf{11.9억} — D+800과 같다 & §1.5 \\
시리즈A & \textbf{80억 클로징} — 리드는 정본에 없다 & §1.5 · §4.13 \\
온담 & 개발 21명 중 \textbf{14명 투입} · 자동차 58사 릴리스 \textbf{7개월 공백} & §1.5 \\
검산 & 14 × 7 = \textbf{98 mm} — 분류 후 86의 114\%, 전부 138의 71\% & day8\_enterprise.py \\
\end{tbl}
\keymsg{상한 6 FTE(연 72 mm)가 계약에 있었다면 나올 수 없는 숫자 — 들어가지 않았거나, 지켜지지 않았거나. 정본은 말하지 않는다}
\note{워크북 §7.1 · 정본 §1.5·§4.13. 세 대조 — ㊱ 항목 5의 예측 vs 98 mm / ㊴ 2×2의 남은 두 칸(한강 / 코너스톤) 중 어느 것이 정합적인가(풀 12\%·이사회·클로징 일정) / "이 결말이 좋은 결말인지는 여러분이 판정한다" — ARR 11.9 그대로 = 여섯 달 순증 0 = 7개월 공백의 값. 결과로 결정을 판정하지 않는다, 결과는 조건을 판정한다. 예상 소요 8분.}
\end{frame}

\begin{frame}{5.7 8일 최종 명제}
\begin{quote}
\large "FN 2.4\%는 자동차 3세트 60건에서 나온 값이다. 온담의 4번째 세트에서는 한 번도 측정되지 않았다. 이것은 여러분이 Day3에 재고 정확도 91.3\%의 분모를 따져 물었을 때 이미 배운 것이다. 그때 여러분은 감사자였다. 지금 여러분은 그 문서를 쓴 사람이다."
\end{quote}
\begin{itemize}
\item 숫자에는 모집단이 있다 — 2.4\%도 91.3\%도
\item 모집단은 옮겨지지 않는다 — 세그먼트를 바꾸면 다시 잰다
\item 자리가 바뀌어도 규칙은 같다 — 감사자로서 따져 물은 것을 서명자로서 적는다
\end{itemize}
\keymsg{모집단 칸을 비운 채 서명한 문서를 센다 — ㉗ v1.0 · ㉞ v2 · ㊲ 초안 · IR 덱 18.9 — 그리고 여러분 회사의 문서}
\note{정본 §9. 조별 발표의 답을 여기서 함께 센다. 예상 소요 6분.}
\end{frame}

\begin{frame}[shrink=8]{학술 배경 — 깊이 읽기 (과정 후)}
\begin{tbl}[\footnotesize]{L{50mm}L{58mm}L{20mm}}
문헌 & 오늘 어디에 & 교재 \\ \midrule
Moore, \textit{Crossing the Chasm} (3판, 2014) & 볼링 핀·전체 제품 — "자동차를 벗어날 것인가" & 6.1 \\
Fisher \& Ury, \textit{Getting to Yes} (3판, 2011) & BATNA(개발하라)·객관적 기준 — 배상 4구간·대한기공 병행 & 6.2 \\
Feld \& Mendelson, \textit{Venture Deals} (4판, 2019) · YC SAFE & economics → control · 옵션풀 셔플 · cap의 정의 & 6.3 \\
Ward \& Daniel, \textit{Benefits Management} (2판, 2012) · MSP · PMBOK 8판 & benefit owner / change owner — 세 서명 & 6.4 \\
ITIL 4 SLM · Beyer 외 SRE (2016) · ISO/IEC 42001:2023 & SLA 3층 · error budget · AI 책임 경계·데이터 대표성 & 6.5 \\
\end{tbl}
\keymsg{네 갈래가 만나는 지점이 ㊵ 항목 4 — 세그먼트(Moore)·객관적 기준(Fisher \& Ury)·투자자의 질문(Feld)·output 열(Ward \& Daniel)}
\note{교재 제6장 6.1\til{}6.5절 요약. 읽는 순서는 6.6절 표. 판본은 인용 전에 확인(제작 노트). 예상 소요 6분.}
\end{frame}

\begin{frame}{읽을 자료 — 필수 · 심화 · 참고}
\begin{itemize}
\item \textbf{필수(과정 전)} — Moore 3\til{}4장 · Fisher \& Ury 4·6장 · Feld \& Mendelson 4\til{}5장 · Ward \& Daniel 3·6장 · 정본 §3.4·§4.10\til{}4.13·§5·§7.4·§7.5·§9 · PM Day4 워크북 ⑯·⑰·⑳
\item \textbf{심화(과정 후)} — YC SAFE 문서·해설 · ITIL 4 SLM practice guide · SRE 3·4장 · MSP 편익 관리 · Skok · a16z "16"·"16 More" · Perri · Cagan · ISO/IEC 42001
\item \textbf{참고} — Day3 ⑫ · Day5 ㉓·㉕ · Day6 ㉖·㉗·㉘·㉚ · Day7 ㉝·㉞·㉟
\item \textbf{읽지 않아도 되는 것} — 엔터프라이즈 세일즈 방법론(파는 법이 아니라 받은 뒤 무엇을 서명하는가) · 밸류에이션 모델링(pre 320은 협상 결과)
\end{itemize}
\keymsg{워크북 X.5·X.6은 오늘 저녁 — 그리고 자기 회사의 문서로 같은 셈을}
\note{교재 6.7절. 예상 소요 3분.}
\end{frame}

\begin{frame}{Product 트랙을 마치며 — 8일의 끝}
\begin{itemize}
\item \textbf{같은 21억이 두 회사에게} — 4.906 / 37.0 / 24.1 / 82.1. 세 번째 숫자를 우리가 먼저 말한다
\item \textbf{47건을 어떻게 판정하는가} — 5차원 표 먼저 · 3분류 · 상한 6 FTE · ②는 우리가 만들지 않는다
\item \textbf{같은 문서를 반대편에서} — SAC → 218문항 · SLA 3층 → 야간 P1 불성립 · 배상 요구 → 하한 6.12와 4구간 · ⑳ → output 열
\item \textbf{표준 문구의 값} — 1.29\%p = 25.4억 · 온담 거절 = 코너스톤 소멸 + pre 242
\item PM 트랙 20종과 Product 트랙 20종은 다른 문서가 아니다 — 같은 문서를 반대편에서 쓴 것. 어느 편에 서든 표의 첫 열은 분모이고 마지막 열은 실명이다
\end{itemize}
\keymsg{검사원의 47분(㉒)에서 시작해 21억(㊱)과 80억(㊴)으로 끝났고, 마지막 문서(㊵)는 첫 문서(①)의 분모로 돌아갔다}
\note{교재 "Product 트랙을 마치며". 8일 과정 종료. 예상 소요 5분.}
\end{frame}
